\section*{NeurIPS Paper Checklist}

\begin{enumerate}

	\item {\bf Claims}
	\item[] Question: Do the main claims made in the abstract and introduction accurately reflect the paper's contributions and scope?
	\item[] Answer: \answerYes{} % Replace by \answerYes{}, \answerNo{}, or \answerNA{}.
	\item[] Justification: The abstract and introduction summarize the results
	      shown in \autoref{fig:single} as discussed in \autoref{sec:perf}. Although
	      they don't directly report all numbers, the headline number is correct, and
	      the introduction states under which conditions it holds. The abstract
	      repeats the same conditions in less precise terms (``up to'', ``large'', ``a
	      given''). Let use note that, however, the correctness of our claims ultimately
	      relies on the correctness on the code in the software package and the benchmark harness:
	      the consistency of the paper is a lesser proxy. A serious assessment would
	      involve checking the unit tests and running the software.

	\item {\bf Limitations}
	\item[] Question: Does the paper discuss the limitations of the work performed by the authors?
	\item[] Answer: \answerYes{} % Replace by \answerYes{}, \answerNo{}, or \answerNA{}.
	\item[] Justification: The most important limitations are discussed in the conclusions (\autoref{sec:conclusions}). More precise but less important considerations of all sorts can be made by perusing \autoref{fig:time-all}.

	\item {\bf Theory assumptions and proofs}
	\item[] Question: For each theoretical result, does the paper provide the full set of assumptions and a complete (and correct) proof?
	\item[] Answer: \answerNA{} % Replace by \answerYes{}, \answerNo{}, or \answerNA{}.
	\item[] Justification: This paper does not really have math. Possibly the DGP in \autoref{eq:dgp} could be considered theory, but its correctness is checked by unit tests in the software package. The considerations in \autoref{sec:gpumem} on the number of trees are more theoretical in nature, but they are a collection of empirical evidence guided by elementary statistical theory without precise mathematical statements.

	\item {\bf Experimental result reproducibility}
	\item[] Question: Does the paper fully disclose all the information needed to reproduce the main experimental results of the paper to the extent that it affects the main claims and/or conclusions of the paper (regardless of whether the code and data are provided or not)?
	\item[] Answer: \answerYes{} % Replace by \answerYes{}, \answerNo{}, or \answerNA{}.
	\item[] Justification: The paper is fully reproducible with a released software package and accompanying scripts to run the benchmarks and produce the figures as included in the paper.

	\item {\bf Open access to data and code}
	\item[] Question: Does the paper provide open access to the data and code, with sufficient instructions to faithfully reproduce the main experimental results, as described in supplemental material?
	\item[] Answer: \answerYes{} % Replace by \answerYes{}, \answerNo{}, or \answerNA{}.
	\item[] Justification: See above.

	\item {\bf Experimental setting/details}
	\item[] Question: Does the paper specify all the training and test details (e.g., data splits, hyperparameters, how they were chosen, type of optimizer) necessary to understand the results?
	\item[] Answer: \answerYes{}
	\item[] Justification: The hyperparameters are mostly left to default
	      values, and all methods are well-known. The DGP for synthetic data is fully
	      specified in \autoref{eq:dgp}.

	\item {\bf Experiment statistical significance}
	\item[] Question: Does the paper report error bars suitably and correctly defined or other appropriate information about the statistical significance of the experiments?
	\item[] Answer: \answerNo{} % Replace by \answerYes{}, \answerNo{}, or \answerNA{}.
	\item[] Justification: To keep the crowded plots easier to parse, and for
	      some benchmarks to reduce the running time, there are no uncertainties.
	      Since all results are repeated measurements across a fine grid, we think
	      that glancing the variability in the trend lines is sufficient to understand
	      the statistical accuracy of the measurements. For what concerns systematic
	      biases, we briefly commented about them in the main text and explained the
	      methodology.

	\item {\bf Experiments compute resources}
	\item[] Question: For each experiment, does the paper provide sufficient information on the computer resources (type of compute workers, memory, time of execution) needed to reproduce the experiments?
	\item[] Answer: \answerYes{} % Replace by \answerYes{}, \answerNo{}, or \answerNA{}.
	\item[] Justification: The CPU and GPU models are clearly reported. The
	      supplementary code includes approximate total running times in the README.
	      Note: though the guidelines suggest to mention the compute that went into
	      unreported experiments, we think it does not make much sense in this
	      context, and we guess it is a concern directed at large neural network
	      training, or more generally computational procedures that output expensive
	      permanent artifacts.

	\item {\bf Code of ethics}
	\item[] Question: Does the research conducted in the paper conform, in every respect, with the NeurIPS Code of Ethics \url{https://neurips.cc/public/EthicsGuidelines}?
	\item[] Answer: \answerYes{} % Replace by \answerYes{}, \answerNo{}, or \answerNA{}.
	\item[] Justification: this work does not have occasions to violate ethical rules.

	\item {\bf Broader impacts}
	\item[] Question: Does the paper discuss both potential positive societal impacts and negative societal impacts of the work performed?
	\item[] Answer: \answerNo{} % Replace by \answerYes{}, \answerNo{}, or \answerNA{}.
	\item[] Justification: The paper introduces no substantially new capabilities.

	\item {\bf Safeguards}
	\item[] Question: Does the paper describe safeguards that have been put in place for responsible release of data or models that have a high risk for misuse (e.g., pre-trained language models, image generators, or scraped datasets)?
	\item[] Answer: \answerNA{} % Replace by \answerYes{}, \answerNo{}, or \answerNA{}.
	\item[] Justification: Nothing more to say.

	\item {\bf Licenses for existing assets}
	\item[] Question: Are the creators or original owners of assets (e.g., code, data, models), used in the paper, properly credited and are the license and terms of use explicitly mentioned and properly respected?
	\item[] Answer: \answerYes{} % Replace by \answerYes{}, \answerNo{}, or \answerNA{}.
	\item[] Justification: We bibliographically cite all the important open
	      source software packages we use.

	\item {\bf New assets}
	\item[] Question: Are new assets introduced in the paper well documented and is the documentation provided alongside the assets?
	\item[] Answer: \answerYes{} % Replace by \answerYes{}, \answerNo{}, or \answerNA{}.
	\item[] Justification: Our released software package is well documented. Even our supplementary material has documentation.

	\item {\bf Crowdsourcing and research with human subjects}
	\item[] Question: For crowdsourcing experiments and research with human subjects, does the paper include the full text of instructions given to participants and screenshots, if applicable, as well as details about compensation (if any)?
	\item[] Answer: \answerNA{} % Replace by \answerYes{}, \answerNo{}, or \answerNA{}.
	\item[] Justification: Nothing more to say.

	\item {\bf Institutional review board (IRB) approvals or equivalent for research with human subjects}
	\item[] Question: Does the paper describe potential risks incurred by study participants, whether such risks were disclosed to the subjects, and whether Institutional Review Board (IRB) approvals (or an equivalent approval/review based on the requirements of your country or institution) were obtained?
	\item[] Answer: \answerNA{} % Replace by \answerYes{}, \answerNo{}, or \answerNA{}.
	\item[] Justification: Nothing more to say.

	\item {\bf Declaration of LLM usage}
	\item[] Question: Does the paper describe the usage of LLMs if it is an important, original, or non-standard component of the core methods in this research? Note that if the LLM is used only for writing, editing, or formatting purposes and does \emph{not} impact the core methodology, scientific rigor, or originality of the research, declaration is not required.
	\item[] Answer: \answerNA{} % Replace by \answerYes{}, \answerNo{}, or \answerNA{}.
	\item[] Justification: Nothing more to say.

\end{enumerate}
